% !TeX root = ../thesis.tex
%
% Drafting scaffold: a staging area for per-paper analysis, not a thesis
% chapter. Write a note as soon as a paper is read; when it finds a home,
% mine it into real prose and delete it. Before submission delete this file
% and its \input in thesis.tex.
%
% One entry per paper, free-form inside (\paragraph headings suit dense
% papers; plain paragraphs are fine too):
%
%   \begin{papernote}{Author Year -- short descriptive title}{BibKey}
%   ...
%   \end{papernote}
%
% BibKey (the references.bib key) becomes \label{note:BibKey}, so any entry
% can be pointed at with \ref{note:BibKey}.

% Defined here, not in thesis.tex, so the scaffold deletes as one file.
\NewDocumentEnvironment{papernote}{mm}
	{\section{#1}\label{note:#2}}
	{}

\chapter{Working Paper Notes}
\label{ch:paper-notes}

\topics{one papernote per paper analysed, structured however the paper calls
    for; keep concrete numbers (resolution, timing, cost figures) verbatim
    with units, since these are what the real chapters will want to cite.}

\begin{papernote}{Thull et al.\ 2026 -- Uncertainty-aware gamma interaction localisation and reconstruction}{Thull2026}

\paragraph{Capsule}
Thull et al.\ \cite{Thull2026} train per-detector, per-axis localisation networks on a
24-block semi-monolithic LYSO preclinical scanner that predict a positional variance
alongside each position estimate, then use that variance to filter and weight LORs in
TOF-OSEM reconstruction --- entirely on measured data. Headline result: the median
LOR-to-point-source distance falls from \mm{20.7} unfiltered to \mm{0.983} under combined
energy and variance filtering at 50\% retention, with variance the stronger and
complementary criterion because energy is a non-monotonic proxy for localisation quality
while predicted variance is monotonic. The work demonstrates rather than characterises:
$\sigma$ is rank-ordered but never calibrated, timing information is discarded, the
detector is fixed, and calibration cost is severe (weeks per scanner, 55 networks).
For this thesis it is the standing evidence that localisation confidence is a
first-class quantity, and it marks the light-distribution route as taken --- leaving the
temporal axis and information-theoretic bounds open.

\paragraph{Summary}
Aachen (Schulz group, RWTH; two authors from Hyperion Hybrid Imaging Systems GmbH)
train localisation networks that emit a per-event positional variance alongside the
position estimate, then use that variance to filter and weight LORs during TOF-OSEM
reconstruction. Fully experimental: no simulation anywhere in the paper. The system is a
proof-of-concept preclinical scanner of 24 semi-monolithic LYSO blocks at \mm{62.5} inner
radius on the Hyperion III platform. Each block is
\qtyproduct{48.4 x 48.3 x 10}{\milli\metre} (planar-segmented,
planar-monolithic, DOI), built from two arrays of 44 polished LYSO slabs of
\qtyproduct{1 x 24 x 10}{\milli\metre} separated by \mm{0.1} BaSO$_4$, with a \mm{0.3}
light guide onto a 12 x 12
channel SiPM array formed from a 6 x 6 grid of Philips DPC3200-22 digital SiPMs.
Baseline detector performance is 12.7\% energy resolution and approximately \ps{500} CTR.

\paragraph{Key Findings}
\begin{itemize}
    \item \textbf{Unfiltered localisation error dominates rather than dilutes.} Median
        distance from reconstructed LOR to a known point-source location is
        \textbf{\mm{20.7}} with no filtering. Energy windowing at 50\% retention brings this
        to \mm{2.4}; variance filtering at 50\% retention to \mm{1.12}; combined energy and
        variance filtering to \textbf{\mm{0.983}}. This is the single most quotable number
        in the paper.
    \item \textbf{Variance outperforms energy as a filtering criterion, and the two are
        complementary.} Mean absolute error at 50\% singles retention (segmented /
        monolithic / DOI, in mm): energy filtering 0.49 / 0.72 / 1.03; variance filtering
        0.25 / 0.53 / 0.90; variance after energy (49\% effective retention)
        0.17 / 0.47 / 0.83. Unfiltered CNN classifier is approximately \mm{1.1} in all three
        dimensions.
    \item \textbf{Energy is a non-monotonic proxy for localisation quality; predicted
        variance is monotonic.} MAE versus deposited energy shows a minimum at \keV{511}, a
        second minimum near \keV{340} (Compton backscatter, \degs{180}, producing a single
        light cluster), and elevated error between them --- a single Compton interaction
        cannot deposit intermediate energy, so that band is necessarily multi-scatter.
        Above \keV{511}, error rises again from pile-up and $^{176}$Lu background. MAE versus
        predicted $\sigma$ is monotonic with growing dispersion, i.e.\ a heavy-tailed error
        distribution.
    \item \textbf{A low-energy, low-variance population exists} (most visible in the
        monolithic and DOI dimensions) that energy windowing discards but variance
        filtering retains; the authors speculate these are geometrically well-constrained
        events occurring directly above a SiPM.
    \item \textbf{Image quality improves, at a noise cost.} With statistics restored to
        250 M coincidences, 50\% singles filtering gives PSNR 14.603 $\to$ 14.899,
        SSIM 0.405 $\to$ 0.462, COV 4.90\% $\to$ 4.66\%. Without restoring counts
        (62.5 M coincidences) COV degrades to 7.49\%. Peak-to-valley ratio on the \mm{0.8}
        rods rises from 1.0985 to 1.1840 under filtering and to 1.2171 under LOR weighting
        at $p = 4$, but COV at $p = 4$ blows out to 10.19\%. Filtering dominated weighting
        at matched effort; recovery coefficients improved at every rod diameter under both.
    \item \textbf{CNN classifiers beat regressors in-plane; CNN regressors win DOI.}
        Classifiers show systematically lower bias, and flood maps show regressors
        producing regions of near-zero prediction density where the continuous output
        collapses.
\end{itemize}

\paragraph{Methods Worth Noting}
\begin{itemize}
    \item \textbf{Heteroscedastic training.} Regressors train with a Gaussian
        negative log-likelihood,
        \[\mathrm{NLL} = \sum_n \bigl[\tfrac{1}{2}\log \sigma_n^2
            + (y_n - \hat{y}_n)^2 / 2\sigma_n^2\bigr],\]
        jointly predicting position and variance. Classifiers use ordinary
        cross-entropy over position bins, with uncertainty taken as the \emph{variance of
        the predicted softmax distribution} --- no special loss required.
    \item \textbf{Per-axis decomposition.} One independent network per detector per spatial
        dimension; 55 networks in total (7 central detectors x 3, 1 central x 2, 16 outer
        x 2). Optuna HPO over 12 configurations, approximately 377 trials each.
    \item \textbf{Fan-beam calibration.} Physical collimated fan of \mm{0.5} width on a
        motorised stage with \um{2} repeatability. Segmented axis: 44 positions at \mm{1.1}
        pitch, 20{,}000 samples each. Monolithic: 96 positions at \mm{0.5} pitch, 9{,}166
        each. DOI: 40 positions at \mm{0.25} pitch, 22{,}000 each. Roughly 880{,}000 events
        per detector per dimension, split 60/20/20. \textbf{Acquisition for the full
        scanner took several weeks.} DOI calibration fires \emph{parallel} to the
        photosensor face, so the lateral coordinate is deliberately unconstrained and
        marginalised over --- the per-axis decomposition is what makes this admissible.
        Only the 8 central-ring detectors could be labelled in all three dimensions; the 16
        outer detectors have monolithic and DOI labels only and borrow a segmented model.
    \item \textbf{Triggering is inherited, not designed.} The paper does not describe a
        trigger scheme; it uses the DPC die-level two-stage trigger and validation (Frach
        2009). The only stated acquisition-time selection is retention of readouts with
        more than 500 detected photons; no energy window is applied until evaluation.
    \item \textbf{Normalisation is recomputed under each filter and weight condition}
        (fan-sum method), which is what keeps the result quantitative rather than
        cosmetically sharper. Uncertainty is combined as
        $\sigma_{3D} = \sqrt{\sigma^2_{\mathrm{seg}} + \sigma^2_{\mathrm{mono}} +
        \sigma^2_{\mathrm{doi}}}$, and LOR weights as
        $w_{\mathrm{LOR}} = 1 / (\sigma_{3D,I_1} \cdot \sigma_{3D,I_2})^p$.
    \item \textbf{Reconstruction:} list-mode TOF-OSEM, 20 iterations, 10 subsets,
        \mm{0.2} isotropic voxels, 250 M coincidences,
        \qty[separate-uncertainty=true]{511 \pm 80}{\kilo\electronvolt} window,
        20 LORs sampled per coincidence. Phantoms: cylindrical shell (normalisation),
        HotRod with 2.0/1.5/1.2/1.0/0.9/\mm{0.8} rods at approximately \MBq{5} initial
        $^{18}$F-FDG, and a \MBq{1.5} $^{22}$Na point source at \mm{25} and \mm{35} radii.
    \item \textbf{Interpretation ceiling.} Label precision is \mm{0.5} (monolithic) and
        \mm{0.25} (DOI) with a \mm{0.5} fan width, so sub-\mm{0.5} MAE figures are not
        interpretable --- the authors say so, and introduce the LOR-to-point-source metric
        specifically to sidestep this.
\end{itemize}

\paragraph{Relevance to This Thesis}
This is the strongest published evidence for the premise that poorly localised
interactions actively damage reconstruction rather than weakly informing it, and it
forecloses the naive form of that idea: per-event uncertainty propagated into
reconstruction via filtering and weighting is done, on real hardware, with normalisation
handled correctly. Cite the \mm{20.7} unfiltered median LOR distance and the non-monotonic
energy--MAE curve as motivation; do not re-derive them.

What the paper does \emph{not} do:
\begin{itemize}
    \item \textbf{It demonstrates rather than characterises.} No Fisher information, no
        CRLB, no bound of any kind. One detector topology, one activity, an empirical
        threshold sweep. There is no crossover analysis for when rejecting a fraction $f$
        buys back more in bias than it costs at $1/\sqrt{1-f}$ in variance.
    \item \textbf{$\sigma$ is rank-ordered, never calibrated.} That $p = 4$ outperforms
        $p = 2$ is a direct tell: if $\sigma$ were a calibrated standard deviation,
        inverse-variance weighting would be optimal by construction. The authors do not
        remark on this. A rank-correct $\sigma$ does not transfer between detectors and
        does not compose across the two singles of a coincidence in any principled way.
    \item \textbf{Timing is entirely unused.} Network inputs are the 144 SiPM photon counts
        and nothing else. The 36 dies supply up to 36 independent timestamps per event at
        approximately \ps{500} CTR, all discarded, along with the pattern of \emph{which}
        dies triggered. The authors state that Compton interaction \emph{sequence} cannot
        be inferred from the light pattern because scintillation emission is isotropic ---
        true of photon counts, but not of arrival times.
    \item \textbf{The detector is fixed.} No design feedback; $\sigma$ is a measurement
        outcome, never an objective.
    \item \textbf{Aleatoric and epistemic uncertainty are not disentangled}, which the
        authors acknowledge: a deterministic network emitting a variance is nominally
        claiming aleatoric uncertainty but is contaminated by finite training data and
        model capacity.
\end{itemize}

\paragraph{Where It Might Go}
Motivation for any chapter arguing that localisation confidence is a first-class quantity;
the \mm{20.7} figure belongs in an introduction. Also the reference point for calibration
burden (several weeks per scanner, 55 networks, outer-ring detectors physically
unlabellable in the segmented direction) --- though note the group's own in-system
calibration work (Kuhl et al.\ 2025) claims to reduce this to hours via virtual
collimation, so ``calibrate faster'' is taken and ``design so that calibration is cheap''
is not.

\emph{Competitive note:} the discussion explicitly announces future work analysing
multi-cluster events, relating light-spread structure to predicted uncertainty, and
characterising the impact of inter-crystal Compton scatter on positioning uncertainty.
Aachen has the built scanner, the fan-beam rig, four RTX 6000 Ada GPUs and a commercial
arm. The light-distribution version of ICS characterisation should be treated as theirs.
The temporal axis --- cardinality and interaction ordering inferred from the joint
distribution of per-channel arrival times rather than from photon counts --- and the
information-theoretic bound on separability are the surviving directions. Note that
Carra et al.\ 2022 runs light-informs-timing and must be distinguished carefully; nobody
appears to run it in reverse.
\end{papernote}